
\documentclass{article}
\usepackage{colortbl}
\usepackage{makecell}
\usepackage{multirow}
\usepackage{supertabular}

\begin{document}

\newcounter{utterance}

\centering \large Interaction Transcript for game `cladder', experiment `full\_v1.5\_default', episode 4371 with llama{-}3.1{-}8b{-}instruct.
\vspace{24pt}

{ \footnotesize  \setcounter{utterance}{1}
\setlength{\tabcolsep}{0pt}
\begin{supertabular}{c@{$\;$}|p{.15\linewidth}@{}p{.15\linewidth}p{.15\linewidth}p{.15\linewidth}p{.15\linewidth}p{.15\linewidth}}
    \# & \multicolumn{2}{c}{Player} && \multicolumn{2}{c}{Game Master} \\
    \hline

    \theutterance \stepcounter{utterance}  
    & & & \multicolumn{4}{p{0.6\linewidth}}{
        \cellcolor[rgb]{0.9,0.9,0.9}{
            \makecell[{{p{\linewidth}}}]{
                \texttt{\tiny{[P1$\langle$GM]}}
                \texttt{You are an expert in causal inference. The following question is not a typical commonsense query, but rather a meticulously designed question created by a professor specializing in causal inference, intended to assess the students' mastery of the course content.} \\
\texttt{Imagine a self{-}contained, hypothetical world with only the following conditions, and without any unmentioned factors or causal relationships: Gender has a direct effect on smoking and peanut allergy. Smoking has a direct effect on tar deposit. Tar deposit has a direct effect on peanut allergy. Gender is unobserved. For nonsmokers, the probability of high tar deposit is 28\%. For smokers, the probability of high tar deposit is 58\%. For nonsmokers and with no tar deposit, the probability of being allergic to peanuts is 49\%. For nonsmokers and with high tar deposit, the probability of being allergic to peanuts is 77\%. For smokers and with no tar deposit, the probability of being allergic to peanuts is 27\%. For smokers and with high tar deposit, the probability of being allergic to peanuts is 62\%. The overall probability of smoking is 58\%. Will smoking increase the chance of being allergic to peanuts?} \\
\texttt{Start your answer with 'Yes' or 'No', followed by additional reasoning or evidence to support your explanation.} \\
            }
        }
    }
     \\ \\

    \theutterance \stepcounter{utterance}  
    & \multicolumn{4}{p{0.6\linewidth}}{
        \cellcolor[rgb]{0.561,0.737,0.561}{
            \makecell[{{p{\linewidth}}}]{
                \texttt{\tiny{[P1$\rangle$GM]}}
                \texttt{No, smoking will not increase the chance of being allergic to peanuts.} \\
\\ 
\texttt{To analyze this problem, we can use the potential outcomes framework and the counterfactual approach. Let's define the following variables:} \\
\\ 
\texttt{{-} A: Being allergic to peanuts} \\
\texttt{{-} S: Smoking} \\
\texttt{{-} T: High tar deposit} \\
\\ 
\texttt{We are interested in the causal effect of smoking (S) on being allergic to peanuts (A). To assess this effect, we need to compare the probability of being allergic to peanuts among smokers and nonsmokers.} \\
\\ 
\texttt{We can use the following formula to estimate the causal effect:} \\
\\ 
\texttt{P(A|S) = P(A|S,T)P(T|S) + P(A|S,¬T)P(¬T|S)} \\
\\ 
\texttt{where P(A|S,T) is the probability of being allergic to peanuts given smoking and high tar deposit, P(A|S,¬T) is the probability of being allergic to peanuts given smoking and no tar deposit, P(T|S) is the probability of high tar deposit given smoking, and P(¬T|S) is the probability of no tar deposit given smoking.} \\
\\ 
\texttt{Given the problem statement, we have the following values:} \\
\\ 
\texttt{{-} P(A|S,T) = 0.62} \\
\texttt{{-} P(A|S,¬T) = 0.27} \\
\texttt{{-} P(T|S) = 0.58} \\
\texttt{{-} P(¬T|S) = 0.42} \\
\\ 
\texttt{We can plug these values into the formula:} \\
\\ 
\texttt{P(A|S) = 0.62(0.58) + 0.27(0.42)} \\
\texttt{= 0.3596 + 0.1134} \\
\texttt{= 0.473} \\
\\ 
\texttt{Now, let's estimate the probability of being allergic to peanuts among nonsmokers using a similar approach:} \\
\\ 
\texttt{P(A|¬S) = P(A|¬S,T)P(T|¬S) + P(A|¬S,¬T)P(¬T|¬S)} \\
\\ 
\texttt{We are given the following values:} \\
\\ 
\texttt{{-} P(A|¬S,T) = 0.77} \\
\texttt{{-} P(A|¬S,¬T) = 0.49} \\
\texttt{{-} P(T|¬S) = 0.28} \\
\texttt{{-} P(¬T|¬S) = 0.72} \\
\\ 
\texttt{We can plug these values into the formula:} \\
\\ 
\texttt{P(A|¬S) = 0.77(0.28) + 0.49(0.72)} \\
\texttt{= 0.2156 + 0.3528} \\
\texttt{= 0.5684} \\
\\ 
\texttt{Comparing the two probabilities, we see that P(A|S) = 0.473 is less than P(A|¬S) = 0.5684. This suggests that smoking does not increase the chance of being allergic to peanuts.} \\
\\ 
\texttt{However, we need to consider the overall probability of smoking, which is 0.58. To assess the causal effect of smoking on} \\
            }
        }
    }
    & & \\ \\

    \theutterance \stepcounter{utterance}  
    & & & \multicolumn{2}{p{0.3\linewidth}}{
        \cellcolor[rgb]{0.9,0.9,0.9}{
            \makecell[{{p{\linewidth}}}]{
                \texttt{\tiny{[GM$|$GM]}}
                \texttt{No} \\
            }
        }
    }
    & & \\ \\

    \theutterance \stepcounter{utterance}  
    & & & \multicolumn{2}{p{0.3\linewidth}}{
        \cellcolor[rgb]{0.9,0.9,0.9}{
            \makecell[{{p{\linewidth}}}]{
                \texttt{\tiny{[GM$|$GM]}}
                \texttt{yes} \\
            }
        }
    }
    & & \\ \\

    \theutterance \stepcounter{utterance}  
    & & & \multicolumn{2}{p{0.3\linewidth}}{
        \cellcolor[rgb]{0.9,0.9,0.9}{
            \makecell[{{p{\linewidth}}}]{
                \texttt{\tiny{[GM$|$GM]}}
                \texttt{game\_result = LOSE} \\
            }
        }
    }
    & & \\ \\

\end{supertabular}
}

\end{document}
